\chapter{Что измеряем: как сравнить модель со шлифом}\label{ch:measures}

Модель выдаёт список пластинок. Опыт --- фотографию шлифа. Чтобы их сравнивать, модель должна
мерить себя \emph{тем же способом}, что и экспериментатор, со всеми особенностями способа.

\section{Доля радиальных гидридов (RHF)}

\fig[0.9]{F21_RHF}{Вес следа гидрида по углу к окружности (Simon; так же в программе PROPHET,
которой пользуются Son и др.): до $40^\circ$ --- 0, от 40 до $65^\circ$ --- 0.5 («смешанные»),
круче $65^\circ$ --- 1. RHF --- взвешенная по длине доля.}

\begin{keyf}[RHF]
\[
\RHF=\frac{\sum_i L_i\,w(\varphi_i)}{\sum_i L_i},\qquad
w(\varphi)=\begin{cases}0,&\varphi\le40^\circ\\0.5,&40^\circ<\varphi<65^\circ\\1,&\varphi\ge65^\circ\end{cases}
\]
\end{keyf}
Важно, \emph{что} считается одним гидридом. На шлифе при увеличении $\times100$--$300$ видна не
отдельная пластинка 0.6~мкм, а пакет или стопка; ориентация пакета может отличаться от ориентации
пластинок в нём. Поэтому в модели два способа:
\begin{itemize}
\item \textbf{по пластинкам} --- прямо по списку (так делать нельзя для сравнения со шлифом, но
  полезно для понимания);
\item \textbf{по изображению} --- пластинки «рисуются» в картинку 0.65~мкм/пиксель с размытием
  1~мкм и шумом, как на снимке СЭМ $\times300$; ориентация измеряется локально на масштабе 3~мкм
  (Фурье-анализ снимка, \code{hydride\_spec.py}). Это и есть «RHF по снимку» в таблицах.
\end{itemize}
Доля длины под $45$--$135^\circ$ к окружности на том же масштабе --- мера Cinbiz ($F_{n45}$).
RHF меньше 0.05 по правилам PROPHET считается шумом.

\begin{exercise}
На шлифе 5 следов: 10~мкм под $10^\circ$, 6~мкм под $50^\circ$, 8~мкм под $80^\circ$, 4~мкм под
$30^\circ$, 2~мкм под $70^\circ$. Найдите RHF. А если 6-микронный след на самом деле стопка из
трёх окружных пластинок по 2~мкм, лежащих ступенькой под $50^\circ$, --- что покажет подсчёт
«по пластинкам»?
\end{exercise}

\section{Доля межзёренных гидридов}

Доля длины пластинок, лежащих на границах зёрен. Son и др. (2026) мерили её по EBSD: в
холоднодеформированном Zr--Nb 92\%, в частично рекристаллизованном 68\%, у радиальных --- 62\%.
В модели --- доля длины пластинок, родившихся на гранях (и их продолжений).

\section{Связность: путь трещины}

Для разрушения важна не только доля, но и то, насколько радиальные гидриды \emph{связаны}.
\begin{itemize}
\item \textbf{RHLD} --- плотность длины радиальных гидридов (мм на мм$^2$).
\item \textbf{HCC} --- доля толщины, перекрытая проекциями гидридов в полосе 100~мкм.
\item \textbf{RHCF} --- самая длинная связная радиальная цепочка в долях толщины.
\item \textbf{RHCP} --- путь трещины наименьшей цены сквозь стенку (рис.~\ref{fig:F22_RHCP}).
\end{itemize}

\fig{F22_RHCP}{RHCP. Шаг по металлу стоит 1, по гидриду --- $1/50$. Алгоритм Дейкстры находит
путь наименьшей цены от внутренней поверхности к наружной; RHCP --- насколько этот путь дешевле
пути по чистому металлу.}

\begin{keyf}[RHCP (как в PROPHET, Kim и др. 2022)]
\[
\mathrm{RHCP}=\frac{1-\text{цена}/H}{1-1/50},
\]
$H$ --- толщина (в клетках). 0 --- на пути нет гидридов, 1 --- путь целиком по гидриду.
\end{keyf}

\paragraph{Зазоры.} Две пластинки, разделённые перемычкой металла в 1~мкм, для трещины --- не
одно и то же, что сплошной гидрид, а на оптическом снимке ($\sim0.5$--$1$~мкм на пиксель) они
сливаются. Поэтому связность считается с разным допуском на зазор: 1~мкм --- «по-настоящему»,
3~мкм --- «как на оптике». В модели радиальные цепочки с зазором оптики --- 30--50~мкм.

\begin{exercise}
Стенка 600~мкм. Трещина идёт 400~мкм по металлу и 200~мкм по гидриду. Чему равен RHCP?
Сколько микрон должно идти по гидриду, чтобы RHCP был 0.5?
\end{exercise}

\section{Двумерное сечение трёхмерной структуры}

Шлиф --- плоское сечение. Диск радиусом $R$, разрезанный плоскостью на расстоянии $d$ от центра,
даёт след длиной $2\sqrt{R^2-d^2}$; наклонный диск даёт след под углом, который зависит от
ориентации плоскости сечения. Поэтому одна и та же трёхмерная структура даёт разные RHF в
поперечном ($r$--$\theta$) и продольном ($r$--$z$) сечениях, а мелкие диски дают много коротких
следов. Двумерная модель сразу живёт в плоскости шлифа и считает пластинки бесконечными по оси
трубы --- это и удобство, и ограничение (гл.~\ref{ch:limits}).

\begin{exercise}
Диски радиусом $R$ разбросаны случайно, сечение --- случайная плоскость. Найдите среднюю длину
следа диска, который сечение задело. (Ответ: $\pi R/2$.)
\end{exercise}

\begin{codebox}
\textbf{Где в коде.} \code{ca\_analysis.py} (\code{simon\_w}, \code{image\_rhf}, пакеты),
\code{calib\_beta.py} (\code{metrics}, \code{image\_measures}), \code{connectivity.py} (RHLD, HCC,
RHCF, RHCP), \code{ca3d\_kinetic.section\_measures} (сечения в 3D).
\end{codebox}
